\section{Software and Hardware Requirements}

\subsection{Software Requirements}

Table~\ref{tab:software} lists the software that each way of running the system requires.
The browser version on the hosted website needs only a web browser.
Building the desktop application and the command line tool requires a Java Development Kit and Apache Maven.

\begin{table}[H]
\centering
\small
\caption{Software requirements}
\label{tab:software}
\begin{tabular}{>{\raggedright\arraybackslash}p{4.2cm}>{\raggedright\arraybackslash}p{3.2cm}>{\raggedright\arraybackslash}p{7.4cm}}
\toprule
\textbf{Software} & \textbf{Version} & \textbf{Purpose} \\
\midrule
Java Development Kit & 21 or newer & Compiles and runs the desktop application and the command line tool \\
Apache Maven & 3 & Builds the project, downloads the libraries, and runs the tests \\
Web browser & Current version & Runs the browser application \\
Git & Any & Downloads the source code, as an alternative to a ZIP file \\
Python & 3 & Serves the browser version locally, which is optional \\
GTKWave & Any & Opens exported VCD waveform files, which is optional \\
\bottomrule
\end{tabular}
\end{table}

Maven downloads the libraries in Table~\ref{tab:libraries} during the first build.
The first build therefore needs an internet connection.
Later builds use the downloaded copies.

\begin{table}[H]
\centering
\small
\caption{Libraries that Maven downloads}
\label{tab:libraries}
\begin{tabular}{llp{8.2cm}}
\toprule
\textbf{Library} & \textbf{Version} & \textbf{Purpose} \\
\midrule
TeaVM & 0.10.2 & Translates the Java engine to JavaScript for the browser version \\
JUnit Jupiter & 5.10.0 & Runs the automated tests \\
JaCoCo & 0.8.14 & Measures the test coverage \\
\bottomrule
\end{tabular}
\end{table}

The system runs on Windows, macOS, and Linux, because both Java and the browser version are platform independent.
The automated build on GitHub Actions uses Ubuntu Linux with Java 21.
The browser checks used the Chromium browser.

\subsection{Hardware Requirements}

The system needs no special hardware.
It runs on an ordinary desktop or laptop computer that can run Java 21 or a current web browser.
It needs no FPGA board, logic analyzer, or other electronic equipment, because it simulates every circuit in software.
A display with a resolution of at least 1280 by 720 pixels shows the desktop and browser interfaces without crowding.

The project did not measure a minimum memory or processor requirement.
For reference, the performance measurements in the final report ran on a virtual machine with two processors at 2.10 GHz.
On that machine, the engine processed more than one million events per second.
